\documentclass[11pt,a4paper]{article}
\usepackage[T1]{fontenc}
\usepackage[utf8]{inputenc}
\usepackage[english]{babel}
\usepackage{amsmath,amssymb}
\usepackage[margin=2cm]{geometry}
\setlength{\parindent}{0pt}
\setlength{\emergencystretch}{3em}
\allowdisplaybreaks

\title{About \texttt{book\_limits.tex}}
\author{}
\date{}
\begin{document}
\maketitle
This note explains what the test data in \texttt{book\_limits.tex} is used for. That file is input for the engine in \texttt{main.c}; it is read by the test \texttt{test\_limits.c} in this folder.

\begin{enumerate}
\item \textbf{Purpose.} The file holds 156 true statements. Each says that a limit has the value printed as its result in a collection of exercises from school and first-year university courses, for example $\lim_{x \to 0} \frac{\sin x - x}{x^{3}} = -\frac{1}{6}$, $\lim_{x \to 0^{+}} x^{x} = 1$, or $\nexists \lim_{x \to 0} \frac{1}{x}$.

\item \textbf{Where it comes from.} Most statements are the proposed statements of \texttt{photos/dati\_esercizi/limiti.jsonl}, whose printed results were confirmed by an independent numerical oracle (values at $x = a \pm 10^{-k}$ at two precisions far apart). The others come from the exercises of that file without a proposed statement: a pair of one-sided limits split in two, or a limit with a parameter taken at the values of the printed cases (for instance $\gamma = 3, 4, 5$ for $\frac{\tan^{2}(2x) - \sin^{2}(2x)}{x^{\gamma}}$), with the value of each case computed by hand. The last 8 were in \texttt{book\_limits\_beyond.tex} until the engine learnt the scale of growth (rules 29k and 29l).

\item \textbf{How it is written.} The collection writes $\log$ for the natural logarithm; here it is written $\ln$, because for this program $\log$ without a base is the logarithm in base $10$, as in school books. Brackets $[\,]$ used as parentheses are written $\left(\,\right)$, and $|x|$ as $\left|x\right|$. In one exercise (zuccher9-2.7) the printed result is wrong, $-\frac{1}{6}$ instead of $-6$: the statement here has the corrected result, and the printed one is in \texttt{false\_limits.tex}.

\item \textbf{What it exercises.} The truncated expansions of chapter 29 at a finite point and at $\pm\infty$ ($x = a \pm t$, $x = \pm\frac{1}{t}$): rational functions with a common factor, roots with fractional exponents, the notable limits, Taylor expansions with deep cancellations, de l'Hôpital's cases, exponential forms $f^{g} = e^{g \ln f}$, the scale $L = \ln\frac{1}{t}$ for $x \ln x$ and $x^{x}$, bounded factors such as $\sin x$ beside an infinite or infinitesimal one, the scale of growth for exponentials and logarithms, a factor $e^{G}$ beside the expansion ($e^{3x}$ against $e^{2x}$, $2^{-x} x^{x}$, $3^{x^{2}} - 9^{x}$, $e^{\sqrt{x}}$ against a logarithm, $e^{\frac{1}{x}}$ on the two sides of $0$), the substitution $x = e^{\frac{1}{t}}$ for $\frac{\ln x}{\ln \ln x}$ and $\frac{\ln x}{\sqrt{1 + 2\ln^{2} x}}$, one-sided limits and limits that do not exist because the one-sided limits differ.

\item \textbf{Expected result.} The line \texttt{PASS: 156 of 156 statements proved, each with a written proof}, and the same statements asked as requests (without their value) answered with a limit that is then proved.

\item \textbf{How to run it.} From the top folder of the repository:
\begin{verbatim}
cc -std=c11 -O2 -Wall -Wextra -o /tmp/test_limits \
    tests/limits/test_limits.c && /tmp/test_limits
\end{verbatim}
\end{enumerate}
\end{document}
